\section{Latency-Aware Evaluation Framework}
\label{sec:framework}

\subsection{Effective Rate}
Let a method $A$ receive the channel $\mathcal{H}(0)=\{\mbf{F},\mbf{g}_k(0),\mbf{h}_{d,k}(0)\}$, return the phase $\bs{\theta}_A$, and need the wall-clock time $\tau_A$ to do so.

\textbf{Definition 1 (Effective rate).} The effective rate of $A$ is its minimum user rate on the channel that exists when it finishes:
\begin{equation}
    R^{\mathrm{eff}}_A = \min_k R_k\big(\bs{\theta}_A;\,\mathcal{H}(\tau_A)\big),
    \label{eq:effective_rate}
\end{equation}
where the precoder is recomputed from $\mathcal{H}(\tau_A)$.

The precoder is a closed-form digital computation that the BS can refresh every frame from pilots, whereas the RIS phase requires the iterative optimization whose latency is measured here. The same split between a fast precoder update and a slower phase update underlies two-timescale RIS designs~\cite{Zhao2021TWC}. Definition~1 therefore lets the phase age and keeps the beam fresh. Section~\ref{sec:res_beam} measures how much this assumption is worth.

\subsection{Gain--Aging Decomposition and the Intrinsic Optimum}
For an iterative method run for $L$ steps, write $r_0(L)=\min_k R_k(\bs{\theta}_L;\mathcal{H}(0))$ for its rate at $t=0$, $\tau(L)$ for its latency and
\begin{align}
    R^{\mathrm{eff}}(L) &= r_0(L) - \Delta(L), \\
    \Delta(L) &= r_0(L) - \min_k R_k\big(\bs{\theta}_L;\mathcal{H}(\tau(L))\big),
\end{align}
so that $\Delta(L)$ is the rate lost to aging during the method's own run time. Additional steps raise $r_0$ with diminishing returns, while $\Delta$ grows with $\tau(L)$. An additional step pays off only while
\begin{equation}
    r_0(L+1)-r_0(L) \;>\; \Delta(L+1)-\Delta(L),
    \label{eq:marginal}
\end{equation}
and the effective rate peaks at the first $L$ where this inequality fails. We call that point the \emph{intrinsic optimum} $L^\star$. It is a property of the optimizer and its step size, the per-step latency and the speed of the channel together, not of the optimizer alone. A slower step or a faster user raises the right-hand side and moves $L^\star$ down. Equation~\eqref{eq:marginal} is a bookkeeping identity rather than a theorem: whether $R^{\mathrm{eff}}$ is unimodal is an empirical question, which we answer in Section~\ref{sec:res_optimum}.

\subsection{Equal Steps versus Equal Milliseconds}
Learned update rules are usually compared with classical optimizers at an equal number of steps. If one learned step costs $c$ times as much as one classical step, an equal-step comparison gives the learned method $c$ times the compute budget. We therefore also compare at equal measured time and use Pareto dominance.

\textbf{Definition 2 (Dominance).} Configuration $B$ dominates configuration $A$ if $\tau_B\le\tau_A$ and $r_{0,B}\ge r_{0,A}$, with one inequality strict, where $r_0$ is the rate at $t=0$.

If $B$ dominates $A$, there is no latency budget at which $A$ is the better choice. A learned configuration that is dominated by some configuration of the classical optimizer is therefore never worth deploying, regardless of how it compares step for step.

\subsection{Normalized Step Cost}
Channel variation depends on how far a user moves in wavelengths, i.e., on $v\tau/\lambda = f_d\tau$. If an optimizer takes $\Delta T$ per step, its latency after $L$ steps in these units is $L\,x$ with
\begin{equation}
    x = \Delta T\, f_d .
    \label{eq:x}
\end{equation}
In our model the users move $v\tau$ metres in time $\tau$ while everything else is fixed, so scaling the step time $\Delta T$ by a factor is exactly equivalent to scaling the speed $v$ by the same factor at a fixed carrier; Section~\ref{sec:res_scaling} varies $\Delta T$. The same $x$ at a different carrier frequency is only approximately equivalent: at 28~GHz the user moves about 9 times fewer metres for the same $x$, which changes path loss, angles and the array's size in wavelengths.

\subsection{Periodic Re-optimization}
Definition~1 evaluates a phase at the instant it is delivered. A deployed controller instead repeats the optimization with period $U$: at each snapshot it spends $T_{\mathrm{p}}$ on pilots for the cascaded channel, during which no data is transmitted, then $L\Delta T$ on optimization, deploys the result, and waits for the next snapshot ($U \ge T_{\mathrm{p}}+L\Delta T$; \emph{back-to-back} operation means $U = T_{\mathrm{p}}+L\Delta T$).

\textbf{Definition 3 (Time-averaged effective rate).} Let $\bs{\theta}(t)$ be the phase deployed at time $t$, and let $\mathcal{P}$ be the set of pilot intervals. The time-averaged effective rate over a window $\mathcal{W}$ is
\begin{equation}
    \bar{R}^{\mathrm{eff}} = \frac{1}{|\mathcal{W}|}\int_{\mathcal{W}\setminus\mathcal{P}} \min_k R_k\big(\bs{\theta}(t);\,\mathcal{H}(t)\big)\,dt .
    \label{eq:time_avg}
\end{equation}
Each update can start from scratch (\emph{cold start}) or from the phase currently deployed (\emph{warm start}). The period trades pilot overhead against the age of the phase in use, and $L$ trades the quality of each update against its delay.
